\begin{quotation}
An upright rod held by a spring seems a textbook exercise, yet the same principle shields gravitational-wave detectors from ground vibration: when the spring almost cancels gravity, the pendulum sways so slowly that seismic noise is filtered out. Monitoring such devices requires their full state, but sensor records are often incomplete. The full--partial reconstruction mapping is a recent machine-learning framework that fills the gaps by exploiting a classic result of nonlinear dynamics: the history of one measured signal encodes the geometry of the whole motion, from which a missing signal can be read off. The framework is tested here on the spring-coupled inverted pendulum, a system simple enough for every prediction to be checked against exact mechanics and rich enough to become chaotic when damped and driven. A mirror symmetry of the pendulum can make the missing signal ambiguous, and regression then averages two possible answers into an impossible one. A normalizing flow, a neural network that returns a full probability distribution, keeps both answers and recovers the chaotic motion when nine tenths of one sensor's readings are missing, whereas physical constraints imposed as soft penalties help with scarce data only where they can reliably tell the mirror states apart.
\end{quotation}
